\documentclass[10pt,a4paper,portrait]{article}

\usepackage[utf8]{inputenc}
\usepackage[a4paper,portrait,margin=0.6in,top=0.52in,bottom=0.52in]{geometry}
\usepackage{tikz}
\usetikzlibrary{shapes,arrows,positioning,calc,fit,backgrounds}
\usepackage{xcolor}
\usepackage{lmodern}
\usepackage{booktabs}
\usepackage{tabularx}
\usepackage{enumitem}
\usepackage{titlesec}
\usepackage{fancyhdr}
\usepackage{microtype}
\usepackage{graphicx}
\usepackage{grffile}

% =========================================================
% Palette & Styling (Matching Lab 2-7 System)
% =========================================================
\definecolor{primary}{RGB}{24, 43, 73}        % Deep clean navy
\definecolor{secondary}{RGB}{14, 116, 144}   % Teal accent
\definecolor{darkslate}{RGB}{30, 41, 59}     % Normal text color
\definecolor{boxbg}{RGB}{248, 250, 252}      % Soft light box
\definecolor{bordergray}{RGB}{226, 232, 240} % Light border

% Hyperref setup with elegant colors
\usepackage[colorlinks=true,linkcolor=primary,urlcolor=secondary,citecolor=primary]{hyperref}

% Terminal Color Scheme: Modern Light Grey Theme
\definecolor{termbg}{RGB}{246, 248, 250}     % Light grey terminal body
\definecolor{termbar}{RGB}{229, 231, 235}    % Medium light grey title bar
\definecolor{termborder}{RGB}{209, 213, 219} % Crisp border line
\definecolor{termtext}{RGB}{31, 41, 55}      % Dark charcoal text
\definecolor{termuser}{RGB}{22, 101, 52}     % Dark green for user
\definecolor{termhost}{RGB}{21, 128, 61}     % Green for host
\definecolor{termpath}{RGB}{29, 78, 216}     % Blue for directory path
\definecolor{termcmd}{RGB}{17, 24, 39}       % Bold black for commands
\definecolor{termout}{RGB}{55, 65, 81}       % Clear dark grey for outputs
\definecolor{termdir}{RGB}{37, 99, 235}      % Blue for directories
\definecolor{termfile}{RGB}{31, 41, 55}      % Dark for regular files
\definecolor{termhidden}{RGB}{107, 114, 128} % Muted grey for hidden files
\definecolor{termsize}{RGB}{180, 83, 9}      % Amber for sizes
\definecolor{termperm}{RGB}{79, 70, 229}     % Indigo for permissions

% =========================================================
% Section Headings
% =========================================================
\titleformat{\section}
  {\color{primary}\large\bfseries\sffamily\raggedright}
  {\thesection.}{0.5em}{}
\titlespacing*{\section}{0pt}{7pt}{2pt}

\titleformat{\subsection}
  {\color{secondary}\normalsize\bfseries\sffamily\raggedright}
  {\thesubsection}{0.5em}{}
\titlespacing*{\subsection}{0pt}{5pt}{2pt}

% =========================================================
% Headers and Footers
% =========================================================
\pagestyle{fancy}
\fancyhf{}
\renewcommand{\headrulewidth}{0.4pt}
\renewcommand{\footrulewidth}{0pt}
\fancyhead[L]{\small\sffamily\color{darkslate} Operating Systems Lab}
\fancyhead[R]{\small\sffamily\color{darkslate} Lab Report: Ubuntu Installation}
\fancyfoot{}

\fancypagestyle{firstpage}{%
  \fancyhf{}
  \renewcommand{\headrulewidth}{0pt}
  \renewcommand{\footrulewidth}{0pt}
}

% =========================================================
% Terminal Window Macro
% =========================================================
\newcommand{\termprompt}[1][\textasciitilde]{%
  \textbf{\textcolor{termuser}{ratul}}\textcolor{termtext}{@}\textbf{\textcolor{termhost}{firefly}}\textcolor{termtext}{:}\textbf{\textcolor{termpath}{#1}}\textcolor{termcmd}{\$}\ %
}

\newcommand{\terminalwindow}[3][\linewidth]{%
\par\vspace{2pt}%
{\centering
\begin{tikzpicture}
  \node[
      text=termtext,
      font=\ttfamily\fontsize{7.2}{9.2}\selectfont,
      align=left,
      text width=\dimexpr#1-0.95cm\relax,
      inner xsep=6pt,
      inner ysep=6pt,
      anchor=north west
  ] (body) at (0.32cm, -0.56cm) {#3};

  \coordinate (topleft) at (0, 0);
  \coordinate (bottomright) at ($(body.south -| #1, 0) + (0, -5pt)$);
  \coordinate (titlecenter) at ($0.5*(topleft) + 0.5*(#1, 0) + (0, -0.28cm)$);

  \begin{scope}[on background layer]
    \fill[termbg, rounded corners=6pt] (topleft) rectangle (bottomright);
    \begin{scope}
      \clip[rounded corners=6pt] (topleft) rectangle (bottomright);
      \fill[termbar] (topleft) rectangle (#1, -0.56cm);
      \draw[termborder, line width=0.5pt] (0, -0.56cm) -- (#1, -0.56cm);
    \end{scope}
    \draw[termborder, line width=0.6pt, rounded corners=6pt] (topleft) rectangle (bottomright);
  \end{scope}

  \fill[red!80!black] (0.36cm, -0.28cm) circle (0.07cm);
  \fill[yellow!80!black] (0.63cm, -0.28cm) circle (0.07cm);
  \fill[green!75!black] (0.90cm, -0.28cm) circle (0.07cm);

  \node[
      text=gray!40!black,
      font=\sffamily\bfseries\fontsize{7.2}{8}\selectfont,
      anchor=center
  ] at (titlecenter) {#2};
\end{tikzpicture}\par}%
\vspace{2pt}%
}

% =========================================================
% Screenshot Box Macro
% =========================================================
\newcommand{\stepfigure}[3][6.7cm]{%
\par\vspace{3pt}%
{\centering
\begin{tikzpicture}
  \node[inner sep=2.5pt, draw=bordergray, line width=0.6pt, rounded corners=5pt, fill=boxbg] {
    \includegraphics[width=0.76\linewidth, height=#1, keepaspectratio]{#2}
  };
\end{tikzpicture}\par}%
\vspace{3pt}%
{\centering\small\sffamily\textbf{\color{primary}#3}\par}%
\vspace{2pt}%
}

\newcommand{\cmd}[1]{\texttt{\textbf{#1}}}

\begin{document}
\raggedbottom
\thispagestyle{firstpage}

% =========================================================
% Title: Only Experiment Title at the Top
% =========================================================
\begin{center}
  {\Large\bfseries\sffamily\color{primary} Experiment Title: Installation and Configuration of Linux Operating System}\\[3pt]
  {\normalsize\sffamily\textbf{Topic:} Step-by-Step Installation of Ubuntu Linux via Bootable USB}
\end{center}
\vspace{-4pt}

% =========================================================
% 1. Objective
% =========================================================
\section{Objective}
The objectives of this experiment are:
\begin{itemize}[noitemsep,topsep=1pt,leftmargin=1.2em]
  \item Understand the full procedure for installing a modern Linux operating system.
  \item Create a bootable installation media using an official Ubuntu ISO and Rufus.
  \item Configure disk partitioning (Root, Swap, EFI) and set up user account security.
  \item Verify successful operating system deployment using terminal-based commands.
\end{itemize}

% =========================================================
% 2. Theoretical Background
% =========================================================
\section{Theoretical Background}
Linux is a free, open-source Unix-like kernel released by Linus Torvalds in 1991. It provides multi-user security, high stability, and modular architecture, powering world servers, supercomputers, cloud infrastructure, and desktop systems.

\subsection{Popular Linux Distributions}
\vspace{-2pt}
\begin{table}[h!]
\centering\footnotesize
\renewcommand{\arraystretch}{1.06}
\begin{tabularx}{\linewidth}{lp{2.4cm}p{4.2cm}X}
\toprule
\textbf{Distribution} & \textbf{Base System} & \textbf{Primary Focus / Use Case} & \textbf{Package Manager} \\
\midrule
\textbf{Ubuntu} & Debian & Beginner-friendly desktop, servers, and cloud & APT (\texttt{deb}) \\
\textbf{Debian} & Independent & Rock-solid stability and open-source standards & APT (\texttt{deb}) \\
\textbf{Fedora} & Independent & Cutting-edge packages and upstream Linux features & DNF / RPM \\
\textbf{Kali Linux} & Debian & Cybersecurity, penetration testing, and forensics & APT (\texttt{deb}) \\
\bottomrule
\end{tabularx}
\end{table}
\vspace{-6pt}

\subsection{Deployment Methods: Bare Metal vs Virtualization}
\vspace{-2pt}
\begin{table}[h!]
\centering\footnotesize
\renewcommand{\arraystretch}{1.06}
\begin{tabularx}{\linewidth}{p{3.2cm}X}
\toprule
\textbf{Method} & \textbf{Simple Explanation} \\
\midrule
\textbf{Bare Metal (Direct)} & Installed directly onto the physical computer drive for maximum hardware performance. \\
\textbf{Dual Boot} & Installed alongside an existing OS (e.g., Windows), letting the user select OS at startup. \\
\textbf{Virtual Machine (VM)} & Runs inside virtualization software (e.g., VirtualBox), providing safe isolated testing. \\
\bottomrule
\end{tabularx}
\end{table}
\vspace{-6pt}

% =========================================================
% 3. Equipment / Software Requirements
% =========================================================
\section{Equipment / Software Requirements}
\vspace{-2pt}
\begin{table}[h!]
\centering\footnotesize
\renewcommand{\arraystretch}{1.08}
\begin{tabularx}{\linewidth}{lXlX}
\toprule
\textbf{Item} & \textbf{Details} & \textbf{Item} & \textbf{Details} \\
\midrule
\textbf{Target Computer} & PC / Laptop (64-bit x86 architecture) & \textbf{RAM Requirement} & Minimum 4 GB recommended \\
\textbf{Storage Space} & Minimum 25 GB free disk space & \textbf{Flash Drive} & 8 GB+ USB Flash Drive \\
\textbf{Target OS ISO} & Ubuntu Desktop (64-bit ISO image) & \textbf{Flashing Tool} & Rufus / BalenaEtcher \\
\bottomrule
\end{tabularx}
\end{table}

% =========================================================
% 4. Procedure / Steps Overview
% =========================================================
\section{Procedure / Steps}
The installation was carried out sequentially through the following 10 steps:
\begin{enumerate}[noitemsep,topsep=1pt,leftmargin=1.4em]
  \item Download the official Ubuntu Desktop ISO disk image.
  \item Write the ISO to a USB flash drive using Rufus to create a bootable drive.
  \item Boot from USB and select ``Try or Install Ubuntu'' from the GRUB menu.
  \item Initialize live environment (Ubuntu splash screen and loading spinner).
  \item Select keyboard layout and choose manual installation for custom partitioning.
  \item Partition the drive: create EFI System Partition and Root (\cmd{/}) partition.
  \item Configure user credentials: enter name, computer name, username, and password.
  \item Review disk partitioning and configuration settings before confirming.
  \item Execute system installation, package copying, and driver setup.
  \item Complete installation, remove the USB drive, and reboot into Ubuntu.
\end{enumerate}

\newpage

% =========================================================
% Step 1 & Step 2
% =========================================================
\subsection{Step 1: Downloading Ubuntu ISO Image}
Visit the official Ubuntu website (\url{https://ubuntu.com/download/desktop}) and download the latest Ubuntu Desktop 64-bit ISO file. This image contains the complete Linux operating system installer and live environment.

\stepfigure{images/downloading_iso.png}{Figure 1: Downloading Ubuntu Desktop ISO from Official Website}

\vfill

\subsection{Step 2: Creating a Bootable USB Flash Drive}
Insert a USB flash drive (minimum 8 GB) into the host computer. Open Rufus, select the downloaded Ubuntu ISO file, choose \textbf{GPT} partition scheme for UEFI systems, and click \textbf{START} to write the ISO image to the flash drive.

\stepfigure{{images/create bootable usb.png}}{Figure 2: Writing ISO Image to USB Flash Drive Using Rufus}

\vfill
\newpage

% =========================================================
% Step 3 & Step 4
% =========================================================
\subsection{Step 3: Selecting ``Try or Install Ubuntu''}
Insert the bootable USB and restart the computer. Press the Boot Menu key (e.g., \textbf{F12}, \textbf{F9}, or \textbf{Esc}) and select the USB drive. The GNU GRUB bootloader appears on screen. Select \textbf{Try or Install Ubuntu} and press \textbf{Enter}.

\stepfigure{{images/try ubuntu.png}}{Figure 3: GRUB Boot Menu -- Selecting Try or Install Ubuntu}

\vfill

\subsection{Step 4: Booting into the Live Environment}
After selecting the installation option, the Ubuntu boot splash screen appears with the official logo and spinner while the operating system initializes the live environment into RAM.

\stepfigure{images/booting_from_usb.png}{Figure 4: Ubuntu Boot Splash Screen Initializing from USB}

\vfill
\newpage

% =========================================================
% Step 5 & Step 6
% =========================================================
\subsection{Step 5: Selecting Manual Installation Type}
Choose system language and keyboard layout. When prompted for installation type, choose \textbf{Manual installation} (or \textbf{Something else}) to maintain full control over disk partitions and filesystem layout.

\stepfigure{{images/select manual installation.png}}{Figure 5: Selecting Manual Installation for Custom Partitioning}

\vfill

\subsection{Step 6: Disk Partitioning}
Allocate drive space for the operating system:
\begin{itemize}[noitemsep,topsep=1pt,leftmargin=1.2em]
  \item \textbf{EFI System Partition:} Size 500 MB (used by UEFI firmware to load the GRUB bootloader).
  \item \textbf{Root Partition (\cmd{/}):} Remaining space formatted as \textbf{Ext4} journaling filesystem.
\end{itemize}

\stepfigure{images/partition.png}{Figure 6: Configuring EFI and Root Disk Partitions}

\vfill
\newpage

% =========================================================
% Step 7 & Step 8
% =========================================================
\subsection{Step 7: Creating User Account and Password}
Set up user identification and system security:
\begin{itemize}[noitemsep,topsep=1pt,leftmargin=1.2em]
  \item Enter your full name and computer hostname.
  \item Choose a username and a strong security password.
  \item Select \textbf{Require my password to log in} for administrative privilege protection.
\end{itemize}

\stepfigure{{images/create account.png}}{Figure 7: Setting Up Username, Computer Hostname, and Password}

\vfill

\subsection{Step 8: Reviewing Installation Configuration}
Before writing changes to disk, the installer presents a summary of selected partitions, mount points, and timezone settings. Review the list and click \textbf{Install} to confirm changes and begin file writing.

\stepfigure{images/review.png}{Figure 8: Reviewing Partition Summary and Confirming Installation}

\vfill
\newpage

% =========================================================
% Step 9 & Step 10
% =========================================================
\subsection{Step 9: System Installation Process}
The installer formats partitions, copies core Linux system files, configures hardware drivers, and sets up localized packages. A progress bar displays real-time setup status.

\stepfigure{{images/installation started.png}}{Figure 9: Automated System Installation and File Extraction}

\vfill

\subsection{Step 10: Completing Installation and System Reboot}
Once file copying completes, an \textbf{``Installation is complete''} notification appears. Click \textbf{Restart Now}. When prompted by the screen, unplug the USB flash drive and press \textbf{Enter} to restart.

\stepfigure{{images/installation finighed.png}}{Figure 10: Installation Finished Notification and Reboot Prompt}

\vfill
\newpage

% =========================================================
% 5. Result / Output Demonstration
% =========================================================
\section{Result / Output Demonstration}
The computer restarts directly into the newly installed Ubuntu Linux system.

\subsection{Desktop Interface Demonstration}
The machine boots cleanly into the GNOME Desktop Environment, showing the top panel, system dock, application drawer, and desktop workspace.

\stepfigure[4.7cm]{{images/post installation.png}}{Figure 11: Ubuntu Desktop Environment Loaded After First Boot}

\vspace{3pt}

\subsection{Terminal-Based System Verification}
To confirm proper operating system installation, kernel functionality, and hardware allocation, we opened the Linux terminal and executed standard verification commands:

\terminalwindow{ratul@firefly: \textasciitilde\ \textemdash\ bash}{
% Command 1: uname -a
\termprompt[\textasciitilde]\textbf{\textcolor{termcmd}{uname -a}}

\textcolor{termout}{Linux firefly 6.8.0-45-generic \#45-Ubuntu SMP PREEMPT\_DYNAMIC Wed Sep 11 15:25:05 UTC 2024 x86\_64 GNU/Linux}

\vspace{0.06cm}
% Command 2: lsb_release -a
\termprompt[\textasciitilde]\textbf{\textcolor{termcmd}{lsb\_release -a}}

\textcolor{termout}{Distributor ID: Ubuntu}
\textcolor{termout}{Description:    Ubuntu 24.04.1 LTS}
\textcolor{termout}{Release:        24.04}
\textcolor{termout}{Codename:       noble}

\vspace{0.06cm}
% Command 3: whoami
\termprompt[\textasciitilde]\textbf{\textcolor{termcmd}{whoami}}

\textcolor{termout}{ratul}

\vspace{0.06cm}
% Command 4: df -h /
\termprompt[\textasciitilde]\textbf{\textcolor{termcmd}{df -h /}}

\textcolor{termout}{Filesystem      Size  Used Avail Use\% Mounted on}
\textcolor{termout}{/dev/nvme0n1p2   59G  8.2G   48G  15\% /}

\vspace{0.10cm}
% Active cursor
\termprompt[\textasciitilde]\textbf{\textcolor{termcmd}{\rule{0pt}{0.9em}\textunderscore}}
}

\vspace{2pt}\noindent\textbf{Explanation of Verification Outputs:}\par\vspace{1pt}
\begin{itemize}[noitemsep,topsep=1pt,leftmargin=1.2em]
  \item \cmd{uname -a}: Confirmed Linux kernel version \texttt{6.8.0-45-generic} on 64-bit architecture.
  \item \cmd{lsb\_release -a}: Verified official distribution name (\texttt{Ubuntu 24.04.1 LTS noble}).
  \item \cmd{whoami}: Confirmed active session under newly created user \texttt{ratul}.
  \item \cmd{df -h /}: Confirmed root filesystem mounted on \texttt{/dev/nvme0n1p2} with 48 GB available space.
\end{itemize}

% =========================================================
% 6. Conclusion
% =========================================================
\section{Conclusion}
Key takeaways from this lab:
\begin{enumerate}[noitemsep,topsep=1pt,leftmargin=1.4em]
  \item Successfully created a bootable USB drive using Rufus and booted into Ubuntu via UEFI.
  \item Configured manual disk partitioning with dedicated EFI and Ext4 Root filesystems.
  \item Created a secure user account and verified system health through terminal commands.
\end{enumerate}

\end{document}
